
\section{Distâncias e Propriedades Geométricas no Espaço de Representação}
\label{sec:representation_geometry}

O pipeline associa a um conjunto finito de objetos a matriz $\mathbf{G}\in\mathbb{S}_+^n$ da Equação \eqref{eq:transformed_gram}. Como $\mathbf{G}$ é semidefinida positiva, existem um espaço de Hilbert $\mathcal{E}$ e vetores $\varphi_1,\dots,\varphi_n\in\mathcal{E}$ tais que
\begin{equation}\label{eq:empirical_feature_realization}
    G_{ij} = \langle \varphi_i, \varphi_j \rangle_{\mathcal{E}},
\end{equation}
e os produtos internos, normas e distâncias entre eles podem ser obtidos diretamente de $\mathbf{G}$, sem que os $\varphi_i$ sejam construídos. Objetos descritos por fontes heterogêneas, ou sem uma métrica previamente especificada, passam assim a admitir uma geometria mensurável. Esta seção é a contrapartida amostral da Subseção \ref{subsec:metric}: a distância e os conjuntos de nível definidos ali sobre $\mathcal{Z}$ são restritos aos objetos analisados e passam a depender apenas de $\mathbf{G}$, e as quantidades usadas nos experimentos (vizinhanças, faixas de nível e energia espectral) são definidas apenas nessa forma. Essa geometria não é uma propriedade intrínseca do domínio original: ela depende do \textit{kernel} e da transformação adotados, e pipelines diferentes podem induzir noções diferentes de proximidade entre os mesmos objetos.

\subsection{Pseudométrica induzida pela matriz de Gram}
\label{subsec:gram_pseudometric}

A distância entre os objetos $i$ e $j$ é $d_{\mathbf{G}}(i,j) = \|\varphi_i-\varphi_j\|_{\mathcal{E}}$, que, pela Equação \eqref{eq:empirical_feature_realization}, é a versão amostral da Equação \eqref{eq:rkhs_distance}:
\begin{equation}\label{eq:gram_distance}
    d_{\mathbf{G}}^2(i,j) = G_{ii} + G_{jj} - 2G_{ij}.
\end{equation}
Pelos mesmos argumentos da Proposição \ref{prop:dist_props}, $d_{\mathbf{G}}$ é uma pseudométrica sobre os objetos, e a matriz de distâncias $\mathbf{D}_{\mathbf{G}} = (d_{\mathbf{G}}(i,j))_{i,j=1}^n$ é simétrica, tem diagonal nula e satisfaz a desigualdade triangular. Como no Comentário \ref{rem:point_sep}, dois objetos distintos podem ter $d_{\mathbf{G}}(i,j) = 0$, o que ocorre se e somente se $\varphi_i = \varphi_j$. Isso não é uma deficiência da construção, mas a indicação de que o pipeline considera os dois objetos indistinguíveis. A relação $i\sim_{\mathbf{G}}j \iff d_{\mathbf{G}}(i,j)=0$ é de equivalência, as suas classes reúnem os objetos de representação coincidente, que não são necessariamente iguais no domínio original, e sobre o quociente $\{1,\dots,n\}/{\sim_{\mathbf{G}}}$ a função $d_{\mathbf{G}}$ é uma métrica. A equivalência é, portanto, relativa ao pipeline $(K,T)$.

\subsection{Geometria relativa da amostra}
\label{subsec:sample_relative_geometry}

A matriz $\mathbf{D}_{\mathbf{G}}$ descreve como as representações se organizam umas em relação às outras, sem acesso às suas coordenadas, e permite analisar proximidade, vizinhanças, isolamento local e agrupamentos.

\subsubsection*{Vizinhanças induzidas}

Para um objeto de referência $i_0$, a vizinhança de raio $\varepsilon\geq 0$ é
\begin{equation}\label{eq:epsilon_neighborhood}
    \mathcal{N}_{\varepsilon}(i_0) = \left\{ j\in\{1,\dots,n\} : d_{\mathbf{G}}(i_0,j) \leq \varepsilon \right\},
\end{equation}
e $\mathcal{N}_k(i_0)$ denota os $k$ objetos mais próximos de $i_0$. Objetos com muitos vizinhos a pequenas distâncias ocupam regiões densas do espaço induzido, e objetos cujos vizinhos permanecem distantes ocupam regiões isoladas. Como $d_{\mathbf{G}}$ é uma pseudométrica, $\mathcal{N}_0(i_0)$ é a classe de equivalência de $i_0$.

\subsubsection*{Conjuntos e faixas de nível}

Fixado $i_0$, a função distância $p_{i_0}(j) = d_{\mathbf{G}}(i_0,j)$ e o conjunto de nível $\Gamma_\rho(i_0) = \{ j : p_{i_0}(j) = \rho \}$ são as restrições amostrais da função $p$ e da Equação \eqref{eq:level_curve}, com $\Gamma_0(i_0)$ igual à classe de equivalência de $i_0$. Como a amostra é finita, $\Gamma_\rho(i_0)$ é vazio para quase todo $\rho$, e distâncias exatamente iguais são raras. Por isso, utilizam-se faixas de nível de espessura $\delta>0$,
\begin{equation}\label{eq:sample_level_band}
    \Gamma_{\rho,\delta}(i_0) = \left\{ j\in\{1,\dots,n\} : \left| d_{\mathbf{G}}(i_0,j)-\rho \right| \leq\delta \right\},
\end{equation}
que agrupam os objetos cujas distâncias à referência pertencem a uma mesma escala. Quando os objetos são posições de um domínio espacial, como as posições de uma imagem, a função distância pode ser avaliada sobre esse domínio, e as faixas assumem a forma de curvas de nível que mostram a organização da representação em torno da referência.

O Algoritmo \ref{alg:geometria} resume o cálculo dessas quantidades a partir da matriz de Gram, sem que as representações $\varphi_i$ precisem ser obtidas.

\begin{algorithm}[htbp]
\caption{Geometria relativa induzida pela matriz de Gram}
\label{alg:geometria}
\begin{algorithmic}[1]
\Require $\mathbf{G} \in \mathbb{S}_+^n$, objeto de referência $i_0$, número de vizinhos $k$, raio $\rho$ e espessura $\delta$
\Ensure $\mathbf{D}_{\mathbf{G}}$, $\mathcal{N}_k(i_0)$ e $\Gamma_{\rho,\delta}(i_0)$
\For{$i, j = 1,\dots,n$}
    \State $(\mathbf{D}_{\mathbf{G}})_{ij} \gets \sqrt{\max\left(G_{ii} + G_{jj} - 2G_{ij},\, 0\right)}$ \Comment{Equação \eqref{eq:gram_distance}}
\EndFor
\State $\mathcal{N}_k(i_0) \gets$ os $k$ objetos mais próximos de $i_0$ segundo $\mathbf{D}_{\mathbf{G}}$
\State $\Gamma_{\rho,\delta}(i_0) \gets \{\, j : |(\mathbf{D}_{\mathbf{G}})_{i_0 j} - \rho| \leq \delta \,\}$
\end{algorithmic}
\end{algorithm}
\subsection{Custo computacional}\label{subsec:custo}

Esta subseção analisa o custo assintótico dos Algoritmos \ref{alg:pipeline}, \ref{alg:geometria} e \ref{alg:energia}, isto é, do método em si, independentemente da instanciação. O custo de obter as fontes, que depende do modelo analisado, é discutido na Subseção \ref{subsec:reprodutibilidade}. Utiliza-se a notação usual \citep{Cormen2009}: $O(g)$ indica um limite superior, $\Omega(g)$ um limite inferior e $\Theta(g)$ um limite justo, que vale como superior e inferior ao mesmo tempo. Sejam $n$ o número de objetos, $d$ o número de fontes, $c_l$ o custo de avaliar o \textit{kernel} $k_l$ sobre um par de objetos, $r$ o número de modos retidos e $k$ o número de vizinhos. Com \textit{kernels} lineares, angulares ou gaussianos sobre representações de dimensão $D_l$, tem-se $c_l = \Theta(D_l)$, e denota-se $D = \sum_l D_l$.

Veja que os três algoritmos são determinísticos e não tomam decisões que dependam dos valores dos dados: para um mesmo $n$ e uma mesma dimensão das fontes, eles executam sempre as mesmas operações. Por esse motivo, o melhor caso, o pior caso e o caso médio coincidem, e o custo de cada algoritmo pode ser descrito por um limite justo $\Theta$. O limite inferior $\Omega$ é utilizado com outro sentido, o do \textit{problema}: ele indica o mínimo que qualquer algoritmo precisa gastar para produzir a mesma saída, e a distância entre ele e o custo do algoritmo mostra onde há espaço para melhorias.

\textbf{Construção da matriz de Gram.} O Algoritmo \ref{alg:pipeline} avalia o \textit{kernel} de cada fonte sobre os $n(n+1)/2$ pares distintos, o que custa $\Theta(n^2 c_l)$ por fonte e $\Theta(n^2 D)$ no total. As transformações por componente utilizadas neste trabalho, a normalização angular $T_a$, o deslocamento por $\mathbf{J}$, o centramento $T_c$ e a normalização em traço $T_{\mathrm{tr}}$, operam entrada a entrada ou por somas de linhas e colunas, e custam $\Theta(n^2)$ cada uma. Veja que o centramento não exige os produtos $H\widetilde{\mathbf{G}}H$, visto que $(H\widetilde{\mathbf{G}}H)_{ij} = \widetilde{G}_{ij} - \bar{g}_i - \bar{g}_j + \bar{g}$, em que $\bar{g}_i$ é a média da linha $i$ e $\bar{g}$ a média de todas as entradas. O produto de Hadamard das $d$ componentes custa $\Theta(d\,n^2)$. Sendo assim, com transformações entrada a entrada, a construção custa $\Theta(n^2 D)$ em tempo e $\Theta(n^2)$ em memória, visto que as fontes podem ser processadas uma de cada vez. As transformações espectrais da Equação \eqref{eq:spectral_transformation}, como $T_w$, exigem uma decomposição espectral e custam $\Theta(n^3)$ \citep{GolubVanLoan2013}. Qualquer algoritmo que produza $\mathbf{G}$ precisa escrever as suas $n(n+1)/2$ entradas e ler as $nD$ coordenadas das fontes, de modo que o problema tem limite inferior $\Omega(n^2 + nD)$. A distância entre esse limite e $\Theta(n^2 D)$ é exatamente o custo do produto $XX^\top$ das representações, que algoritmos rápidos de multiplicação de matrizes reduzem apenas assintoticamente.

\textbf{Geometria.} No Algoritmo \ref{alg:geometria}, cada entrada da matriz de distâncias é obtida da Equação \eqref{eq:gram_distance} em tempo constante, e a matriz inteira custa $\Theta(n^2)$. Como a saída tem $n^2$ entradas, esse custo coincide com o limite inferior $\Omega(n^2)$, isto é, a construção da geometria é ótima. A vizinhança $\mathcal{N}_k(i_0)$ e a faixa de nível $\Gamma_{\rho,\delta}(i_0)$ de uma referência exigem apenas a linha $i_0$ da matriz e custam $\Theta(n)$, a primeira por seleção em tempo linear \citep{Cormen2009}, e ordenar a linha, como no algoritmo, custa $\Theta(n \log n)$. As leituras que tomam cada objeto como referência, como os $k$-vizinhos \textit{leave-one-out}, custam $\Theta(n^2)$ por seleção e $\Theta(n^2 \log n)$ por ordenação.

\textbf{Energia espectral.} O Algoritmo \ref{alg:energia} realiza a decomposição espectral completa de $\mathbf{G}$, que custa $\Theta(n^3)$ com os algoritmos densos usuais para matrizes simétricas \citep{GolubVanLoan2013}, e a partir dela obtém as energias e a razão $\widehat{E}_{\mathrm{ratio}}$ de todos os objetos em $\Theta(nr)$. Essa é a etapa mais cara do método, mas ela pode ser evitada em duas partes. Primeiro, a razão de participação da Equação \eqref{eq:r_participation} não exige os autovalores, visto que, para uma matriz simétrica, $\sum_k \hat\lambda_k = \operatorname{tr}(\mathbf{G})$ e $\sum_k \hat\lambda_k^2 = \|\mathbf{G}\|_F^2$, de modo que $r$ é obtido em $\Theta(n^2)$. Segundo, as leituras utilizam apenas os $r$ modos retidos, e estes podem ser obtidos por métodos de Lanczos ou aleatorizados em $O(n^2 r)$ \citep{Halko2011}. O problema tem limite inferior $\Omega(n^2)$, visto que $\mathbf{G}$ precisa ser lida.

A Tabela \ref{tab:custo} resume a análise. Com a decomposição espectral completa, o método custa $\Theta(n^2 D + n^3)$ por matriz de Gram, e com a decomposição truncada custa $O\bigl(n^2 (D + r)\bigr)$, contra um limite inferior de $\Omega(n^2 + nD)$. Veja que, em ambos os casos, o custo é polinomial em $n$ e linear na dimensão das fontes, e não depende do número de atributos da entrada nem do número de parâmetros do modelo, que entram apenas na obtenção das fontes.

\begin{table}[pos=htbp]
\centering
\caption{Custo assintótico das etapas do método, para $n$ objetos, $d$ fontes de dimensão total $D$, $r$ modos retidos e $k$ vizinhos. A segunda coluna é o custo dos algoritmos como descritos, a terceira é uma alternativa mais barata, quando existe, e a quarta é o limite inferior do problema. As células vazias indicam que não há alternativa mais barata.}
\label{tab:custo}
\small
\begin{tabular}{l c c c}
\hline
Etapa & Algoritmo & Alternativa & Problema \\
\hline
Gram bruta de todas as fontes & $\Theta(n^2 D)$ & & $\Omega(n^2 + nD)$ \\
$T_a$, deslocamento, $T_c$, $T_{\mathrm{tr}}$, Hadamard & $\Theta(d\,n^2)$ & & $\Omega(n^2)$ \\
Matriz de distâncias & $\Theta(n^2)$ & & $\Omega(n^2)$ \\
$\mathcal{N}_k(i_0)$ e $\Gamma_{\rho,\delta}(i_0)$ & $\Theta(n \log n)$ & $\Theta(n)$ & $\Omega(n)$ \\
$k$-vizinhos de todos os objetos & $\Theta(n^2 \log n)$ & $\Theta(n^2)$ & $\Omega(n^2)$ \\
Razão de participação & $\Theta(n^3)$ & $\Theta(n^2)$ & $\Omega(n^2)$ \\
Decomposição espectral & $\Theta(n^3)$ & $O(n^2 r)$ & $\Omega(n^2)$ \\
Energias e $\widehat{E}_{\mathrm{ratio}}$, dada a decomposição & $\Theta(nr)$ & & $\Omega(nr)$ \\
\hline
Total & $\Theta(n^2 D + n^3)$ & $O\bigl(n^2(D + r)\bigr)$ & $\Omega(n^2 + nD)$ \\
\hline
\end{tabular}
\end{table}